\chapter{Materiales y componentes empleados}
\label{anexo:materiales}

Este anexo recoge los principales componentes utilizados en el prototipo TierraViva. La finalidad no es presentar una lista comercial de compra, sino documentar los elementos físicos empleados, su función dentro del sistema y su estado de validación durante el desarrollo.

Los componentes se han agrupado según su papel dentro de la arquitectura: estación de campo, comunicaciones, sensorización, actuación, alimentación y centro de control. Esta organización facilita relacionar cada material con los bloques descritos en el cuerpo principal de la memoria.

\begin{table}[H]
\begin{center}
\small
\begin{tabular}{|p{0.22\textwidth}|p{0.08\textwidth}|p{0.26\textwidth}|p{0.18\textwidth}|p{0.16\textwidth}|}
\hline
\textbf{Componente} & \textbf{Cant.} & \textbf{Función} & \textbf{Estado} & \textbf{Observaciones} \\
\hline
Arduino UNO & 2 & Controlador de estación de campo. & Probado & Lectura de sensores y control local del relé. \\
\hline
A7670E-LASA & 2 & Comunicación LTE de estación. & Validado por comandos AT & Envío remoto de telemetría y consulta de comandos. \\
\hline
BME280 & 2 & Medición de temperatura, humedad relativa y presión atmosférica. & Probado & Sensor ambiental digital mediante I2C. \\
\hline
BH1750 / GY-302 & 2 & Medición de luminosidad ambiente. & Probado & Sensor digital de luz mediante I2C. \\
\hline
Sensor capacitivo de humedad de suelo V1.2 & 2 & Medición de humedad del sustrato. & Probado parcialmente & Lectura analógica; requiere calibración con suelo seco y húmedo. \\
\hline
Módulo relé & 2 & Conmutación de la alimentación de la bomba. & Pendiente/probado según fase & Actuación local desde Arduino. \\
\hline
Bomba de agua 12 V & 2 & Actuación física sobre el riego. & Pendiente/probado según fase & Estado seguro: bomba apagada. \\
\hline
LM2596 & 2 & Regulación de tensión de 12 V a 5 V. & Pendiente/probado según fase & Debe ajustarse a 5,0 V antes de alimentar la electrónica. \\
\hline
Fuente de alimentación 12 V & 2 & Alimentación principal de estación y bomba. & Pendiente/probado según fase & Base para pruebas de estación en banco. \\
\hline
Raspberry Pi 5 & 1 & Centro de control del sistema. & Probada & Ejecución de SQLite, FastAPI, dashboard y reglas de decisión. \\
\hline
Tarjeta SIM de datos & 1--2 & Acceso a red móvil para el módulo LTE. & Probada parcialmente & Necesaria para comunicación remota de las estaciones. \\
\hline
Resistencias 1 k$\Omega$ y 2 k$\Omega$ & Varias & Divisor de tensión para la línea RXD del A7670E. & Pendiente/probado según fase & Adaptación de nivel desde Arduino hacia el módulo LTE. \\
\hline
Protoboard y cableado Dupont & Varias unidades & Montaje y pruebas de conexiones. & Usado en pruebas & Permite modificar el prototipo durante la validación. \\
\hline
\end{tabular}
\caption{Resumen de componentes empleados en TierraViva}
\label{cuadro:materiales_tierraviva}
\end{center}
\end{table}

La Tabla~\ref{cuadro:materiales_tierraviva} recoge los materiales principales del prototipo. El estado indicado no debe interpretarse como una certificación industrial del componente, sino como el grado de validación alcanzado durante el desarrollo del TFG. En algunos casos, como los sensores o el módulo LTE, se han realizado pruebas específicas de funcionamiento antes de su integración completa. En otros, como la bomba de agua o la alimentación final de estación, la validación queda asociada a las fases de actuación y pruebas de integración.

El uso de componentes habituales en prototipado responde a dos criterios. Por un lado, permite reducir el coste y facilitar la sustitución de piezas durante el desarrollo. Por otro, favorece la reproducibilidad del sistema, ya que se trata de módulos disponibles comercialmente y con documentación accesible. Esta decisión es coherente con el alcance académico del proyecto, cuyo objetivo es validar una arquitectura funcional y no desarrollar una solución industrial cerrada.

\section{Referencias comerciales de compra}
\label{anexo:referencias_compra}

Algunos componentes se adquirieron a través de distribuidores generalistas, principalmente para facilitar la disponibilidad y reducir el tiempo de aprovisionamiento. Estas referencias se incluyen únicamente como trazabilidad de compra. Para la justificación técnica se han utilizado, siempre que ha sido posible, hojas de datos o documentación del fabricante del componente.

\begin{table}[H]
\begin{center}
\small
\begin{tabular}{|p{0.28\textwidth}|p{0.30\textwidth}|p{0.30\textwidth}|}
\hline
\textbf{Componente} & \textbf{Referencia comercial} & \textbf{Uso en el prototipo} \\
\hline
Sensor capacitivo de humedad de suelo & AZDelivery, módulo higrómetro capacitivo V1.2 compatible con Arduino. & Lectura analógica de humedad del sustrato. \\
\hline
BME280 & Módulo BME280 compatible con alimentación de 5 V. & Medición ambiental: temperatura, humedad relativa y presión. \\
\hline
BH1750 / GY-302 & Módulo GY-302 basado en sensor BH1750. & Medición de luminosidad ambiente en lux. \\
\hline
\end{tabular}
\caption{Referencias comerciales de sensores empleados}
\label{cuadro:referencias_comerciales_sensores}
\end{center}
\end{table}

La referencia comercial permite identificar el tipo de módulo utilizado, pero no sustituye a la documentación técnica. Por este motivo, en el cuerpo principal se citan fuentes técnicas específicas para justificar las características principales de cada sensor: documentación de AZ-Delivery para el sensor capacitivo, hoja de datos de Bosch Sensortec para el BME280 y hoja de datos del BH1750FVI para el sensor de luminosidad.